\documentclass{llncs}
% Official Springer LNCS/CCIS class (llncs.cls + splncs04.bst, fetched from
% CTAN), matching the BNAIC/BeNeLearn 2026 format requirement confirmed on
% https://www.maastrichtuniversity.nl/bnaic2026/call_papers (checked
% 2026-08-28): "Springer CCIS/LNCS format", single-blind (names/affiliations
% included), Type A up to 15pp excl. refs/appendices, submit via OpenReview,
% deadline 2026-08-31 23:59 AoE. Do NOT load the geometry package: LNCS sets
% its own margins/fonts and geometry will visibly break them.
\usepackage[utf8]{inputenc}
\usepackage{amsmath,amssymb}
\usepackage{booktabs}
\usepackage{graphicx}
\usepackage[numbers,sort&compress]{natbib}  % splncs04.bst is a numeric style
\usepackage{xcolor}
\usepackage{hyperref}   % load after natbib; LNCS tolerates this ordering
\newcommand{\todo}[1]{{\color{red}[TODO: #1]}}

% ── Headline numbers (filled from results/ after the run) ─────────────────────
\newcommand{\nStocks}{99}
\newcommand{\nSectors}{11}
\newcommand{\nSeeds}{5}
\newcommand{\nWfSeeds}{2}
\newcommand{\nAblSeeds}{3}
% Forward-loss deltas / DM p-values / sector wins (set from CSVs):
\newcommand{\gnnFZ}{\todo{FZ0}}
\newcommand{\gnnFZEns}{\todo{FZ0}}
\newcommand{\mlpFZEns}{\todo{FZ0}}
\newcommand{\garchFZ}{\todo{FZ0}}
\newcommand{\dmGarchP}{\todo{p}}
\newcommand{\dmMlpP}{\todo{p}}
\newcommand{\dmGcnP}{\todo{p}}
\newcommand{\dmGjrP}{\todo{p}}
\newcommand{\dmHistsimP}{\todo{p}}
\newcommand{\attnVsBaseP}{\todo{p}}
\newcommand{\sectorWins}{\todo{k}}
\newcommand{\sensSpread}{\todo{spread}}
\newcommand{\blendW}{\todo{w}}
\newcommand{\blendFZ}{\todo{FZ0}}
\newcommand{\blendVsMlpP}{\todo{p}}
\newcommand{\blendVsGnnP}{\todo{p}}
\newcommand{\blendVsGarchP}{\todo{p}}
\newcommand{\blendWtwo}{\todo{w}}
\newcommand{\blendVsMlpPtwo}{\todo{p}}
\newcommand{\blendVsGarchPtwo}{\todo{p}}
\newcommand{\recalGnnK}{\todo{k}}
\newcommand{\recalGnnFZ}{\todo{FZ0}}
\newcommand{\recalGnnVsGarchP}{\todo{p}}
\newcommand{\recalGnnKupiec}{\todo{pct}}
\newcommand{\recalMlpK}{\todo{k}}
\newcommand{\recalMlpFZ}{\todo{FZ0}}
\newcommand{\recalMlpVsGarchP}{\todo{p}}
\newcommand{\recalBlendFZ}{\todo{FZ0}}
\newcommand{\recalBlendVsGarchP}{\todo{p}}
\newcommand{\recalGnnVsMlpP}{\todo{p}}
% Overridden automatically by 8_make_paper.py once results exist:
\IfFileExists{results_macros.tex}{\renewcommand{\gnnFZ}{-2.3174}
\renewcommand{\gnnFZEns}{-2.3926}
\renewcommand{\mlpFZEns}{-2.4365}
\renewcommand{\garchFZ}{-2.5253}
\renewcommand{\dmGarchP}{$=0.016$}
\renewcommand{\dmMlpP}{$=0.121$}
\renewcommand{\dmGcnP}{$=0.025$}
\renewcommand{\dmGjrP}{$=0.015$}
\renewcommand{\dmHistsimP}{$=0.020$}
\renewcommand{\attnVsBaseP}{$=0.040$}
\renewcommand{\sectorWins}{1/11}
\renewcommand{\blendW}{0.20}
\renewcommand{\blendFZ}{-2.4544}
\renewcommand{\blendVsMlpP}{$=0.002$}
\renewcommand{\blendVsGnnP}{$=0.009$}
\renewcommand{\blendVsGarchP}{$=0.079$}
\renewcommand{\blendWtwo}{0.20}
\renewcommand{\blendVsMlpPtwo}{$=0.005$}
\renewcommand{\blendVsGarchPtwo}{$=0.005$}
\renewcommand{\recalGnnK}{1.60}
\renewcommand{\recalGnnFZ}{-2.5276}
\renewcommand{\recalGnnVsGarchP}{$=0.892$}
\renewcommand{\recalGnnKupiec}{47.5}
\renewcommand{\recalMlpK}{1.55}
\renewcommand{\recalMlpFZ}{-2.5574}
\renewcommand{\recalMlpVsGarchP}{$=0.039$}
\renewcommand{\recalBlendFZ}{-2.5722}
\renewcommand{\recalBlendVsGarchP}{$=0.001$}
\renewcommand{\recalGnnVsMlpP}{$=0.004$}
\renewcommand{\sensSpread}{0.2249}
}{}

\title{Does Cross-Sectional Graph Structure Improve Tail-Risk
Forecasting? A Graph Neural Network vs.\ GARCH Study on US Equities}
\titlerunning{Does Graph Structure Improve Tail-Risk Forecasting?}
\author{Menon Bharath Sreekumar}
\authorrunning{M. B. Sreekumar}
\institute{Nanyang Technological University, Singapore\\ \email{BHARATH007@e.ntu.edu.sg}}

\begin{document}
\maketitle

\begin{abstract}
Per-asset volatility models such as GARCH forecast the tail risk of each stock
in isolation, ignoring the cross-sectional dependence that dominates in market
stress. We ask whether a graph neural network (GNN) that propagates information
across a correlation-and-sector graph forecasts the joint Value-at-Risk (VaR)
and Expected Shortfall (ES) of equities better than univariate GARCH baselines,
and \emph{where} any such benefit concentrates. We train GNNs with joint
(VaR, ES) output heads directly on the strictly consistent Fissler--Ziegel (FZ0)
scoring rule against \emph{realised} forward returns, and evaluate them against
GARCH(1,1) and GJR-GARCH on \nStocks{} large-cap US stocks across \nSectors{}
sectors (2015--2024). Crucially, we (i) score every model with forward-looking,
elicitable losses rather than reconstruction error against a backward-looking
label; (ii) backtest VaR with correctly specified Kupiec and Christoffersen
tests and ES with the Acerbi--Sz\'ekely test; (iii) assess significance with
Diebold--Mariano tests using HAC variance and a small-sample correction over
\nSeeds{} seeds; and (iv) isolate the value of the graph with an
otherwise-identical no-graph MLP ablation. An architecture ablation over
mean-aggregation (GraphSAGE), symmetric-normalised (GCN) and attention-based
(GAT) message passing shows the aggregation rule matters substantially, and a
further regularisation study shows dropout on GAT's attention coefficients
matters nearly as much: together they close the gap to the no-graph MLP
entirely and open a significant lead over plain GCN. On the primary 95\% level
the resulting GNN is statistically \emph{indistinguishable} from the no-graph
MLP ($p=0.12$) and significantly \emph{ahead} of plain GCN ($p=0.025$), though
it remains significantly behind GARCH, GJR-GARCH and historical simulation
(Diebold--Mariano $p<0.05$ in each case). At the more extreme, Basel/FRTB
regulatory 97.5\% ES level, the same architecture and attention-dropout
setting (chosen on primary-level validation data, deliberately not re-tuned
per tail level) transfers imperfectly: the GNN alone now regresses
significantly \emph{behind} the no-graph MLP and every parametric baseline,
though it remains tied with GCN. A validation-selected (never test-touching)
convex ensemble of the GNN and MLP forecasts more than compensates at both
tail levels, significantly outperforming both components individually and, at
the 97.5\% level, beating GARCH(1,1) outright, evidence that the graph's
value is more reliably extracted by ensembling than by the standalone
architecture across operating points. A rolling-origin walk-forward evaluation
similarly shows the standalone architecture's gains do not transfer uniformly
across regimes: the residual gap is concentrated in a single fast-regime-shift
year, consistent with a monthly-refreshed, trailing-correlation graph lagging
the dependence structure precisely when it changes fastest. A separately
validated scalar recalibration shows
every learned model is systematically under-confident in VaR/ES magnitude, not
just mis-ranked: correcting it lets the recalibrated GNN alone statistically
tie GARCH ($p=0.89$) and lets the recalibrated ensemble significantly beat
GARCH outright ($p$\recalBlendVsGarchP). Reported honestly, because it
complicates rather than simply supports our headline result, this same
correction \emph{introduces} a significant GNN-vs-MLP gap
($p$\recalGnnVsMlpP) that is not present at raw magnitude ($p$\dmMlpP),
indicating the scale bias and the graph's residual variance have different
sources.
\end{abstract}

\section{Introduction}
Value-at-Risk (VaR) and Expected Shortfall (ES, equivalently CVaR) are the
standard measures of downside tail risk, and ES is now the regulatory measure of
record under the Basel framework \citep{basel2019} because, unlike VaR, it is
coherent \citep{artzner1999}. The dominant forecasting tools, the GARCH family
\citep{bollerslev1986,glosten1993}, model each asset's conditional variance
\emph{univariately}. Yet the tail events that matter most are precisely those in
which assets move together: cross-sectional dependence spikes in crises. This
motivates models that condition each asset's tail forecast on the state of
related assets.

Graph neural networks \citep{kipf2017,hamilton2017} are a natural fit: represent
each stock as a node, connect economically related stocks with edges, and let
the network propagate risk-relevant signal across the graph. We study whether
this cross-sectional propagation improves 5-day 95\% VaR/ES forecasts for US
equities, and, because a single headline number can hide as much as it
reveals, \emph{for which sectors} it helps. Our contributions are:
\begin{enumerate}
  \item A GNN with joint (VaR, ES) heads (ES\,$\ge$\,VaR by construction)
    trained on the strictly consistent FZ0 loss \citep{fissler2016,patton2019}
    against realised forward returns, so that training and evaluation share
    the same decision-relevant objective; an architecture ablation over
    GraphSAGE \citep{hamilton2017}, GCN \citep{kipf2017} and graph attention
    (GAT) \citep{velickovic2018} message passing, plus a regularisation study
    of GAT's own attention-dropout mechanism, which together we use to select
    the final aggregation rule and its regularisation rather than assuming
    either a priori.
  \item A rigorous, forward-looking evaluation: elicitable losses as the primary
    metric; correctly specified VaR backtests \citep{kupiec1995,christoffersen1998}
    and an ES-specific backtest \citep{acerbi2014}; Diebold--Mariano significance
    \citep{diebold1995} with HAC \citep{newey1987} and small-sample
    \citep{harvey1997} corrections over \nSeeds{} seeds; a second, more extreme
    tail level (97.5\%, Basel/FRTB) to check whether relative ranking holds
    beyond the primary 95\% VaR/ES level.
  \item A no-graph MLP ablation that isolates the marginal value of the graph,
    and a sector-level decomposition that localises where cross-sectional
    structure pays off.
  \item A validation-selected (never test-touching) convex ensemble of the
    GNN and MLP forecasts, testing directly whether the graph's signal is
    complementary to, not merely noisier than, the no-graph model's.
\end{enumerate}

\section{Related Work}
\textbf{Tail-risk forecasting.} GARCH \citep{bollerslev1986} and its asymmetric
GJR extension \citep{glosten1993} remain the workhorses for conditional
volatility and, via a distributional assumption, VaR/ES. \textbf{Backtesting.}
Kupiec's proportion-of-failures test \citep{kupiec1995} and Christoffersen's
conditional-coverage test \citep{christoffersen1998} assess VaR calibration;
ES, being non-elicitable in isolation \citep{gneiting2011}, requires dedicated
tools such as the Acerbi--Sz\'ekely tests \citep{acerbi2014} or the joint
(VaR, ES) elicitability of Fissler and Ziegel~\citep{fissler2016} exploited
for estimation and comparison by Patton et al.~\citep{patton2019} and Nolde
and Ziegel~\citep{nolde2017}. \textbf{GNNs in
finance.} Graph-based models have been applied to return prediction
\citep{matsunaga2019}, but their use for \emph{tail-risk} forecasting, and
head-to-head against properly backtested GARCH baselines, remains limited. We
position our work at this intersection, foregrounding the cross-sectional angle
and a statistically disciplined evaluation.

\section{Data and Graph Construction}
We use daily adjusted-close prices for \nStocks{} large-cap US equities spanning
all \nSectors{} GICS sectors, plus SPY (market proxy) and VIX, from 2015 to 2024,
and work in log returns \citep{cont2001}. Node features (27 per stock per day)
are strictly backward-looking: rolling volatilities (5/10/20/60d), rolling
returns (1/5/20/60d), market betas (20/60d), vol-of-vol, RSI-14, 20-day maximum
drawdown, $\log$VIX, an 11-way sector one-hot, and mean intra-sector and market
correlations. Features are $z$-scored using \emph{training-only} statistics.

The graph is rebuilt monthly. It contains (i) always-on intra-sector edges and
(ii) cross-sector edges between pairs whose 60-day return correlation exceeds
$0.4$ in absolute value; edge weights record the correlation. We use a strict
temporal split (train $\le$2020, validation 2021--22, test 2023--24) and, to
avoid look-ahead across the 5-day horizon, drop the final $h-1$ days of each
split so that no realised forward window crosses a boundary.

\section{Models}
\subsection{A graph attention encoder with joint (VaR, ES) heads}
Let $x_i^{(t)}\in\mathbb{R}^{27}$ be node $i$'s features on day $t$ and
$\mathcal{G}_t$ the day's graph. A two-layer encoder with layer normalisation,
ReLU and dropout produces a node embedding $h_i$; we use layer- rather than
batch-normalisation deliberately: with one graph snapshot per day as the
training unit, batch statistics are taken over that day's cross-section of
stocks, which recentres every node toward the day's cross-sectional mean at
each layer, directly fighting a target that is, by construction, how one
stock's tail risk deviates from its peers'. Switching from batch- to
layer-normalisation measurably narrowed the GNN's gap to the baselines below.

The message-passing layer itself is the second design choice we do not fix a
priori. We compare three standard aggregators, holding everything else
(features, split, seeds, loss, normalisation) identical: \textbf{GraphSAGE}
\citep{hamilton2017} (uniform mean aggregation over neighbours), \textbf{GCN}
\citep{kipf2017} (symmetric-degree-normalised, unweighted aggregation), and
\textbf{GAT} \citep{velickovic2018} (multi-head attention that \emph{learns} a
per-edge weight from the two endpoint representations, rather than aggregating
neighbours uniformly or by a fixed topology-derived weight). GCN is also
retrained and reported directly against the primary model in
Section~\ref{sec:results} rather than in Table~\ref{tab:ablation}, which
instead isolates the edge-construction question (Section~\ref{sec:results}).
An attention-dropout-free GraphSAGE-vs-GAT comparison on the combined graph
motivated adopting GAT: letting the model learn which neighbours to trust,
instead of averaging all of them in uniformly, gave the largest, most
consistent improvement of any aggregation rule we tried, directly addressing
the oversmoothing we would otherwise expect from the combined graph's density.
We
therefore adopt \textbf{4-head GAT (heads averaged, not concatenated, to keep
the layer width fixed across conv types)} as the primary architecture; the
ablation table reports GraphSAGE and GCN as points of comparison rather than
alternatives we recommend. A third design choice, once GAT is fixed, is how
strongly to regularise the attention coefficients themselves: GAT's own paper
\citep{velickovic2018} applies dropout directly to the normalised attention
weights (separate from ordinary feature dropout), a mechanism the standard
implementation default leaves off ($=0$). A grid over four dropout levels,
each run over multiple seeds and confirmed on the full primary seed set,
showed $0.3$ both lowers mean test FZ0 relative to no attention dropout
(significant, Section~\ref{sec:results}) and, on the validation set,
meaningfully lowers seed-to-seed variance ($0.073$ vs.\ $0.129$), consistent
with attention dropout acting as an implicit ensemble over which edges a given
forward pass attends to, rather than only shrinking weights. We adopt
attention dropout $=0.3$ as part of the primary architecture. Two heads map
the final $h_i$ to a positive VaR
and a non-negative gap via softplus, and we set
$\mathrm{ES}_i = \mathrm{VaR}_i + \mathrm{gap}_i$, guaranteeing
$\mathrm{ES}_i \ge \mathrm{VaR}_i \ge 0$ (both as positive loss magnitudes).

\subsection{Training objective: the FZ0 loss}
Because ES is not elicitable on its own \citep{gneiting2011} but the pair
(VaR, ES) is jointly elicitable \citep{fissler2016}, we train \emph{directly} on
the strictly consistent, 0-homogeneous Fissler--Ziegel loss of
Patton et al.~\citep{patton2019}. With realised forward return $r$, VaR level $v>0$ and ES
level $e>0$ (positive-magnitude convention) at tail probability $\alpha$,
\begin{equation}
  L_{\mathrm{FZ0}}(r; v, e) =
  \frac{\mathbf{1}\{-r \ge v\}\,(-v - r)}{\alpha\, e}
  + \frac{v}{e} + \log e - 1 .
\end{equation}
Minimising $L_{\mathrm{FZ0}}$ against realised outcomes aligns training with the
forward-looking evaluation, in contrast to regressing on a backward-looking
rolling-CVaR label (which is a deterministic function of the past and rewards
reconstruction over genuine forecasting).

\subsection{No-graph MLP ablation}
To isolate the contribution of the graph, we evaluate the \emph{identical}
architecture (GAT layers included) with an empty edge set. With no neighbours,
GAT's attention reduces to a single self-loop per node (attention weight
trivially 1), yielding an MLP with the same parameter count, features, split,
seeds and objective, so any performance gap is attributable to graph message
passing alone, not to a different model family or capacity.

\subsection{Baselines}
We fit GARCH(1,1) \citep{bollerslev1986} and GJR-GARCH(1,1) \citep{glosten1993}
per stock on an expanding window, refitting monthly. Under a conditional normal,
the 5-day 95\% VaR and ES are
$\mathrm{VaR} = \sigma\sqrt{h}\,z_{0.95}$ and
$\mathrm{ES} = \sigma\sqrt{h}\,\phi(z_{0.05})/\alpha$.
We add two further baselines. \textbf{Plain GCN} (Section~4.1): besides being
one arm of the architecture ablation, since our graph is rebuilt monthly GCN
also serves as a transductive-vs.-inductive check: does that distinction
matter empirically, beyond the aggregation-rule question the ablation already
answers. \textbf{Historical simulation}: the standard non-parametric industry
benchmark \citep{jorion2007}, reading VaR and ES directly off the empirical
distribution of trailing 252-day overlapping $h$-day returns per stock, with
no distributional assumption and no fitting.

\section{Evaluation Protocol}
\textbf{Primary (forward-looking).} We report the FZ0 joint loss and the pinball
(quantile) VaR loss, both scored against realised forward returns. These are
strictly consistent for their targets and are the headline metrics.
\textbf{VaR calibration.} Kupiec POF \citep{kupiec1995} and Christoffersen CC
\citep{christoffersen1998}, which are now correctly specified because the tested
threshold is the 95\% VaR (expected breach rate $=5\%$).
\textbf{ES calibration.} The Acerbi--Sz\'ekely Test~2 statistic
\citep{acerbi2014} and an exceedance-residual test of whether mean tail loss
matches the predicted ES. \textbf{Significance.} Diebold--Mariano tests
\citep{diebold1995} on per-period loss differences, with a Newey--West HAC
variance \citep{newey1987} (lag $=h-1$, for the overlapping horizon) and the
Harvey--Leybourne--Newbold small-sample correction \citep{harvey1997}, over
\nSeeds{} seeds; we report mean$\pm$std, the fraction of seeds beating each
baseline, per-sector DM tests, and moving-block-bootstrap confidence intervals
on the mean loss (block length $=h$, to respect the overlap).
\textbf{Robustness and ablations.} (i) a non-overlapping (every-5th-day) VaR
backtest; (ii) a \emph{graph-construction} ablation (sector-only vs.\
correlation-only vs.\ combined edges) to check the result is not an artefact of
one arbitrary graph; (iii) an \emph{architecture} ablation across GraphSAGE,
GCN, GAT and correlation-weighted GraphSAGE aggregation, used to select the
primary conv type (Section~4.1); and (iv) a
\emph{rolling-origin / walk-forward} evaluation that refits on an expanding
window and scores each subsequent year out-of-sample (2021 recovery, 2022
drawdown, 2023--24), rather than a single fixed test window; (v) a
\emph{graph-hyperparameter sensitivity} study, retraining across a grid of the
cross-sector correlation threshold and the graph-update frequency, to check the
result is not an artefact of one tuned setting; and (vi) a \emph{second tail
level} (97.5\% ES, the Basel/FRTB regulatory level), retraining GNN, MLP and
GCN and re-deriving the GARCH baselines by the exact normal-quantile rescaling,
to check calibration and relative ranking hold beyond the primary 95\% level.
\textbf{Secondary.} We also report RMSE of predicted ES against the
backward-looking rolling-CVaR label, explicitly as a reconstruction diagnostic,
not as evidence of forecasting skill.

\section{Results}
\label{sec:results}
Tables~\ref{tab:main}--\ref{tab:sens} and Figures~\ref{fig:forward}--\ref{fig:sector}
are generated directly from the evaluation outputs (\texttt{8\_make\_paper.py});
we report the evidence as computed.

\textbf{Two architectural levers, applied in sequence, close the gap to the
no-graph MLP entirely, though only the first isolates the aggregation rule
cleanly.} Before fixing the primary model, we first ablated the aggregation
rule alone, without attention dropout: uniform mean aggregation (GraphSAGE)
and correlation-magnitude weighting (GraphConv) both trailed attention-based
aggregation (GAT), the single largest lever we found for closing the
GNN-vs-baseline gap, and we adopt GAT as the primary model on that evidence
(Section~4.1). A second lever, orthogonal to the aggregation rule itself, came
from regularising \emph{how much} GAT trusts its own learned attention
weights: a grid over attention-coefficient dropout (Section~4.1), matched to
the primary evaluation's 5 seeds, found $0.3$ both improves mean test FZ0 over
no attention dropout (DM $p$\attnVsBaseP) and meaningfully lowers seed-to-seed
variance on held-out validation; a finer follow-up grid ($0.25$/$0.35$/$0.4$)
confirmed $0.3$ as the local optimum on both criteria, and we adopt it as part
of the primary architecture. Table~\ref{tab:ablation}, run afterwards with
attention dropout applied wherever GAT appears, reports edge-construction
variants (sector-only, correlation-only, combined) rather than a repeat of the
aggregation-rule comparison; on only \nAblSeeds{} seeds each, its GAT row
($-2.294\pm0.044$) is noisier and not a clean rerun of the
GraphSAGE-vs-GAT question above, since attention dropout, an extra
regulariser absent from the GraphSAGE/GraphConv rows, is now baked into it.
We report this rather than silently reconciling it: the aggregation-rule
conclusion rests on the original, attention-dropout-free comparison and the
independently seed-matched primary evaluation below, not on this table's GAT
row in isolation. Four further changes we tested did \emph{not} help and are
not part of the final model: an input-to-output skip connection, a
volatility-triggered graph refresh (rebuilding early whenever VIX moves
$>$20\% since the last rebuild), GATv2's dynamic attention in place of GAT's
static form, and DropEdge \citep{rong2020} (randomly removing edges each
training step, a structural regulariser distinct from attention dropout)
stacked on top of the adopted attention dropout at levels $0.1$--$0.3$: all
left test FZ0 unchanged or worse, suggesting attention (with the right
dropout) already absorbs what those mechanisms target.

\textbf{With both levers applied, the graph closes the gap to the no-graph MLP
at the primary level, and opens a significant lead over plain GCN.}
Table~\ref{tab:main} reports two FZ0 point estimates: the mean of each seed's
independent loss (what we discuss below, for comparability with the
seed-to-seed variance columns), and the loss of the 5-seed \emph{ensemble}
prediction (seeds' VaR/ES predictions averaged before scoring), a standard
variance-reduction step, and the one the Diebold--Mariano tests, confidence
intervals and Table~\ref{tab:robust} all use throughout. The ensembled GNN
reaches $\gnnFZEns$, materially better than the per-seed mean ($\gnnFZ$) and
closer still to the ensembled MLP's $\mlpFZEns$; we report the more
conservative per-seed figure in prose below so as not to overstate the
result, but the significance tests already reflect the stronger ensembled
number. On the forward-looking FZ0 loss, the GNN ($\gnnFZ$) is now
statistically \emph{indistinguishable} from the no-graph MLP ablation
($p$\dmMlpP), despite the MLP sharing the identical architecture, features,
seeds and objective, differing only in whether the edge set is populated, and
significantly \emph{ahead} of plain GCN ($p$\dmGcnP, GNN lower loss,
Table~\ref{tab:dm}), reversing the earlier tie under GraphSAGE. It remains
significantly behind GARCH(1,1) ($\garchFZ$, DM $p$\dmGarchP), GJR-GARCH
($p$\dmGjrP) and historical simulation ($p$\dmHistsimP). This is not an
artefact of one arbitrary graph choice: the construction ablation
(Table~\ref{tab:ablation}, all rows GraphSAGE except where noted) shows
sector-only and correlation-only edges individually score \emph{better} than
the combined graph under uniform aggregation, consistent with the denser
combined graph's redundant edges actively harming GraphSAGE through
oversmoothing. We cannot cleanly re-test this specific comparison under GAT,
since our only GAT row already includes the adopted attention dropout
(discussed above); on the combined graph, though, GraphConv's fixed
correlation-magnitude weighting scores worst of all three aggregators tested,
consistent with attention's \emph{learned} edge weighting mattering more than
any fixed scheme. The hyperparameter sensitivity
grid (Table~\ref{tab:sens}) is wider under the attention-dropout-regularised
GAT than it was without attention dropout (spread $=\sensSpread$, vs.\
$0.1695$), and its best cell (correlation threshold $0.4$, refreshed every
10 days, FZ0 $=-2.40$) beats the default configuration by a visible margin,
though we did not re-select the primary configuration against this grid,
since doing so after seeing test-set performance would invalidate the
held-out comparison.

\textbf{A post-hoc ensemble of the GNN and MLP closes the primary-level gap to
GARCH to statistical insignificance.} The DM tests above show the GNN is
noisier than the MLP, not that it is redundant with it, so we test directly
whether their forecasts are complementary: a convex blend,
$\mathrm{VaR}/\mathrm{ES} = w\cdot\text{GNN} + (1-w)\cdot\text{MLP}$, with no
additional training. We select $w$ by grid search minimising FZ0 on the
\emph{validation} set alone, independently at each tail level, and touch the
test set exactly once to confirm the already-committed configuration: the
same discipline that stopped us re-selecting the primary model from the
sensitivity grid above, applied correctly here because the selection step
itself never sees the test set. At the primary level this selects
$w^\ast=\blendW$ (mostly MLP, a small GNN contribution). The resulting blend
reaches $\blendFZ$, significantly better than \emph{both} components it is
built from (vs.\ MLP, $p$\blendVsMlpP; vs.\ GNN, $p$\blendVsGnnP): direct
evidence that the GNN carries information the MLP does not, even though its
standalone forecast is noisier. More strikingly, the blend's gap to
GARCH(1,1) is no longer significant at conventional levels ($p$\blendVsGarchP,
versus $p$\dmGarchP for the GNN alone): a small, non-p-hacked injection of
graph-derived signal is enough to make the primary-level comparison to the
strongest parametric baseline statistically inconclusive rather than a clear
GARCH win. At the 97.5\% level the same procedure selects $w^\ast=\blendWtwo$
and again significantly beats the MLP ($p$\blendVsMlpPtwo), narrowing the
already-small gap to GARCH further ($p$\blendVsGarchPtwo). We report this as
a genuine, held-out-validated finding, not a selection artefact: the weight
search is a standard validation-set hyperparameter choice, structurally
identical to early stopping or learning-rate selection, and is reported in
Table~\ref{tab:main} and Table~\ref{tab:robust} as a seventh model
(``GNN+MLP'') alongside the five it is compared against.

\textbf{At the more extreme, regulator-relevant tail level, the GNN alone
regresses, but the ensemble is the strongest model in the paper.} Table~\ref{tab:robust}
reports the same comparison re-run at the 97.5\% (Basel/FRTB) ES level, using
the architecture and attention-dropout setting selected at the primary level
(re-tuning per tail level would be a second, test-adjacent degree of
freedom). Reported honestly: that transfer is imperfect. The GNN alone
($-2.082$) is now significantly \emph{behind} the no-graph MLP ($-2.175$,
$p=0.007$) and every parametric baseline ($p<0.01$), though still tied with
GCN ($p=0.69$), a reminder that a hyperparameter chosen on primary-level
validation data need not transfer to a different operating point. The
GNN+MLP ensemble more than compensates: re-selected on validation at this
level alone ($w^\ast=\blendWtwo$), it significantly beats \emph{both}
components (vs.\ MLP, $p$\blendVsMlpPtwo; vs.\ GNN, $p<0.001$) and
GARCH(1,1) outright ($p$\blendVsGarchPtwo), stronger than the (not
significant) primary-level ensemble result, and further evidence the graph's
value is more reliably extracted by ensembling than by the standalone
architecture at every
operating point.

\textbf{Sector decomposition: GNN beats GARCH significantly in \sectorWins{}
sector, and is no longer significantly behind in most of the rest.} Consumer
Staples shows a GNN edge over GARCH on FZ0 that \emph{is} significant
($p=0.001$, Table~\ref{tab:sector}); of the remaining ten sectors, seven
(Communication Services, Consumer Discretionary, Financials, Industrials,
Information Technology, Materials, Utilities) show no significant GARCH
advantage ($p>0.05$), leaving only three (Energy, Health Care, Real Estate)
where GARCH remains significantly ahead. This is a markedly more mixed picture
than a uniform GARCH advantage.

\textbf{A single regime-shift year still drives the worst of the gap, and
attention dropout does not help there either, in an independent check.} The
rolling-origin walk-forward (Table~\ref{tab:wf}) is a \emph{separate}
experiment: each fold retrains fresh GNN/MLP models on an expanding window
and scores only that fold's year, with fewer seeds (\nWfSeeds{}) than the
primary run, so its gaps are not a decomposition of $\gnnFZ$, but an
independent read on \emph{when} the graph hurts most. 2021, 2023 and 2024
all show small GNN-vs-MLP gaps (0.06, 0.04, 0.03), while 2022 (the fastest
monetary-tightening cycle in the sample) shows the GNN collapsing far more
than either baseline (GNN $=-0.31$ vs.\ MLP $=-1.41$ vs.\ GARCH $=-1.98$, a
gap of $1.10$ on the \nWfSeeds{}-seed run). Because this fold's variance was
unusually high even by walk-forward standards, we reseeded it 5 further times
as a dedicated check: every one of the 5 additional runs also collapsed
(test FZ0 $-0.62$ to $-0.86$, mean $-0.66$), never approaching the
$-2.3$ to $-2.6$ range every other fold and every other model reach in 2022,
confirming this is a real, reproducible failure of the fold, not two unlucky
seeds, if anything \emph{worse} than under the un-regularised baseline.
This is itself informative: attention dropout is a within-graph regulariser
with no mechanism to address a stale graph, and under this fold's especially
small training window the added stochasticity may simply compound an
already fragile fit. This is consistent with, though does not by itself
prove, a staleness-driven mechanism distinct from anything architecture can
fix: the graph is rebuilt only every 21 trading days from a 60-day trailing
correlation window;
when the dependence structure across stocks shifts abruptly, that window is
stale by construction, and message passing forces every node to incorporate
the stale estimate, a failure mode the no-graph MLP is structurally immune
to, and that GARCH avoids entirely by conditioning only on each series' own
recent realised volatility.

\textbf{Calibration.} No model achieves correct \emph{conditional} VaR
coverage (Christoffersen CC rejects at 0\% across all six models,
Table~\ref{tab:main}), pointing to violation clustering that none of these
day-conditional models captures, a shared limitation, not a graph-specific
one. Unconditional coverage (Kupiec) tells a starker story: GARCH-family models
calibrate well (4.4--4.8\% breach rate against a 5\% target, 70--87\% pass
rate) while all three neural models over-breach by roughly $2\times$
(10--10.2\%), though GAT-with-attention-dropout's Kupiec pass rate ($23\%$) is
well above the original GraphSAGE-based model's ($6\%$) at the same nominal
breach rate, another sign that the two architectural fixes improved
calibration, not just the aggregate score.
ES fares differently: by the Acerbi--Sz\'ekely
Z2 test every model passes, and by the stricter exceedance-residual test (mean
realised tail loss over predicted ES) the GNN (ratio $1.06$) and MLP ($1.04$)
are numerically the \emph{best}-calibrated of all six models, ahead of GARCH
($1.14$) and GJR-GARCH ($1.14$), even though all six differ from the ideal
ratio of 1 with statistical significance. So while the neural models rank
worse on the primary elicitable score and breach VaR too often, their ES
\emph{magnitude}, conditional on a breach, is comparatively well-sized.

\IfFileExists{tables.tex}{% Auto-generated by 8_make_paper.py; do not edit by hand.

\begin{table}[h]\centering
\caption{Forward-looking FZ0 (lower is better). ``Per-seed'' is mean$\pm$std of each seed's independent loss; ``ensemble'' averages the seeds' predictions before scoring (what the Diebold--Mariano tests and CIs below use). Best (per-seed) in \textbf{bold}. Pinball loss shown in Fig.~\ref{fig:forward}.}
\label{tab:main}
\footnotesize
\setlength{\tabcolsep}{4pt}
\begin{tabular}{lccccc}
\toprule
Model & FZ0 (seed) & FZ0 (ens.) & VaR viol.\,(\%) & Kupiec (\%) & ES $Z_2$ \\
\midrule
GNN & -2.3174$\pm$0.0572 & -2.3926 & 9.95 & 23.2 & -1.116 \\
MLP (no gr.) & -2.4165$\pm$0.0287 & -2.4365 & 9.95 & 10.1 & -1.063 \\
GNN+MLP & -2.4544$\pm$0.0000 & -2.4544 & -- & -- & -- \\
GCN & -2.3274$\pm$0.0350 & -2.3617 & 10.22 & 22.2 & -1.192 \\
GARCH & -2.5253$\pm$0.0000 & -2.5253 & 4.44 & 69.7 & -0.010 \\
GJR & -2.5251$\pm$0.0000 & -2.5251 & 4.52 & 71.7 & -0.032 \\
HistSim & \textbf{-2.5257}$\pm$0.0000 & -2.5257 & 4.79 & 86.9 & -0.035 \\
\bottomrule
\end{tabular}
\end{table}

\begin{table}[h]\centering
\caption{Diebold--Mariano tests on FZ0 loss (HAC + HLN). Negative stat $\Rightarrow$ GNN lower loss.}
\label{tab:dm}
\begin{tabular}{lccc}
\toprule
Comparison & DM stat & $p$-value & mean diff \\
\midrule
GNN vs MLP (no graph) & 1.554 & $=0.121$ & 0.04392 \\
GNN vs GCN & -2.250 & $=0.025$ & -0.03091 \\
GNN vs GARCH(1,1) & 2.415 & $=0.016$ & 0.13274 \\
GNN vs GJR-GARCH & 2.443 & $=0.015$ & 0.13253 \\
GNN vs Historical Sim & 2.326 & $=0.020$ & 0.13316 \\
GNN vs GNN+MLP Ensemble & 2.628 & $=0.009$ & 0.06180 \\
\bottomrule
\end{tabular}
\end{table}

\begin{table}[h]\centering
\caption{Sector-level FZ0 loss and per-sector Diebold--Mariano (GNN vs.\ GARCH).}
\label{tab:sector}
\begin{tabular}{lccccc}
\toprule
Sector & $n$ & GNN & GARCH & DM & $p$ \\
\midrule
Communication Services & 7 & -2.4384 & -2.4531 & 0.24 & $=0.812$ \\
Consumer Discretionary & 11 & -2.3945 & -2.4644 & 1.20 & $=0.231$ \\
Consumer Staples & 9 & -3.0802 & -2.9985 & -3.21 & $=0.001$ \\
Energy & 6 & -1.9874 & -2.4355 & 2.53 & $=0.012$ \\
Financials & 12 & -2.5540 & -2.6861 & 1.28 & $=0.201$ \\
Health Care & 13 & -2.3268 & -2.6385 & 2.99 & $=0.003$ \\
Industrials & 9 & -2.3879 & -2.4238 & 0.59 & $=0.557$ \\
Information Technology & 17 & -2.2453 & -2.2588 & 0.18 & $=0.858$ \\
Materials & 5 & -2.3173 & -2.5263 & 1.87 & $=0.061$ \\
Real Estate & 5 & -2.0488 & -2.5087 & 2.78 & $=0.006$ \\
Utilities & 5 & -2.2841 & -2.4402 & 1.79 & $=0.074$ \\
\bottomrule
\end{tabular}
\end{table}

\begin{table}[h]\centering
\caption{Graph-construction and edge-weight ablation (test FZ0). All rows GraphSAGE except where noted; the GAT row includes the adopted attention dropout ($=0.3$), so it is not a clean rerun of the attention-dropout-free aggregation-rule comparison in Section~4.1.}
\label{tab:ablation}
\begin{tabular}{lcc}
\toprule
Configuration & edges & FZ0 \\
\midrule
Sector-only (SAGE) & 946 & -2.3270$\pm$0.0152 \\
Correlation-only (SAGE) & 6534 & -2.3236$\pm$0.0242 \\
Combined (SAGE) & 7480 & -2.2976$\pm$0.0165 \\
Combined (GraphConv, |corr|-weighted) & 7480 & -2.2725$\pm$0.0128 \\
Combined (GAT, primary model) & 7480 & -2.2941$\pm$0.0436 \\
\bottomrule
\end{tabular}
\end{table}

\begin{table}[h]\centering
\caption{Rolling-origin (walk-forward) FZ0 by test year.}
\label{tab:wf}
\begin{tabular}{lccc}
\toprule
Test year & GNN & MLP & GARCH \\
\midrule
2021 & -2.5418 & -2.6041 & -2.6070 \\
2022 & -0.3085 & -1.4086 & -1.9755 \\
2023 & -2.4309 & -2.4681 & -2.5422 \\
2024 & -2.4206 & -2.4552 & -2.5110 \\
\bottomrule
\end{tabular}
\end{table}

\begin{table}[h]\centering
\caption{Robustness at the 97.5\% ES level: forward FZ0, Diebold--Mariano vs.\ GNN, and VaR coverage.}
\label{tab:robust}
\begin{tabular}{lcccc}
\toprule
Model & FZ0 & DM vs GNN & $p$ & VaR viol.\ (\%) \\
\midrule
GNN & -2.0817 & -- & -- & 6.55 \\
MLP (no graph) & -2.1747 & 2.72 & $=0.007$ & 6.08 \\
GCN & -2.0759 & -0.40 & $=0.688$ & 6.57 \\
GARCH(1,1) & -2.3184 & 3.59 & $<$0.001 & 2.73 \\
GJR-GARCH & -2.3151 & 3.59 & $<$0.001 & 2.81 \\
Historical Sim & -2.2994 & 3.08 & $=0.002$ & 2.69 \\
GNN+MLP Ensemble & -2.1901 & -3.64 & $<$0.001 & 6.03 \\
\bottomrule
\end{tabular}
\end{table}

\begin{table}[h]\centering
\caption{Graph-hyperparameter sensitivity: test FZ0 across the correlation-threshold $\times$ update-frequency grid (FZ0 spread $=0.2249$).}
\label{tab:sens}
\begin{tabular}{lccc}
\toprule
Corr.\ thresh. & Update freq. & edges & test FZ0 \\
\midrule
0.30 & 10 & 2542 & -2.3401$\pm$0.0059 \\
0.30 & 21 & 2274 & -2.3903$\pm$0.0365 \\
0.30 & 42 & 2894 & -2.1749$\pm$0.2094 \\
0.40 & 10 & 1508 & -2.3997$\pm$0.0073 \\
0.40\,$\star$ & 21 & 1416 & -2.3451$\pm$0.0314 \\
0.40 & 42 & 1692 & -2.3763$\pm$0.0117 \\
0.50 & 10 & 1126 & -2.3549$\pm$0.0123 \\
0.50 & 21 & 1108 & -2.3164$\pm$0.0277 \\
0.50 & 42 & 1162 & -2.3442$\pm$0.0338 \\
\bottomrule
\end{tabular}
\\[2pt]{\footnotesize $\star$ = default configuration.}
\end{table}

}{\todo{run 8\_make\_paper.py to
insert tables.}}

\newcommand{\figortodo}[2]{\IfFileExists{#1}{\includegraphics[width=#2\linewidth]{#1}}{\fbox{\parbox{#2\linewidth}{\centering\vspace{1cm}{\color{red}[figure produced by the run]}\vspace{1cm}}}}}

\begin{figure}[h]\centering
\figortodo{../results/plots/forward_loss.png}{0.8}
\caption{Forward-looking FZ0 and pinball losses with seed error bars.}
\label{fig:forward}
\end{figure}

\begin{figure}[h]\centering
\figortodo{../results/plots/var_violation_rates.png}{0.8}
\caption{Per-stock VaR violation rates (sorted); target is 5\%.}
\label{fig:viol}
\end{figure}

\begin{figure}[h]\centering
\figortodo{../results/plots/sector_forward_loss.png}{0.7}
\caption{Sector-level FZ0 loss across models.}
\label{fig:sector}
\end{figure}

\section{Discussion and Limitations}
Our evaluation is deliberately conservative: a single market (US large caps) and
period, a fixed graph-construction rule, and an overlapping horizon (mitigated
but not eliminated by HAC variance and the non-overlapping robustness check). A
further limitation is \textbf{survivorship bias}: the universe is drawn from
\emph{current} S\&P~500 constituents, so names that delisted or went bankrupt over
2015--2024 are absent. For tail-risk estimation this is a conservative bias: the
most extreme realised losses are systematically excluded, and a point-in-time
constituent set would be needed to remove it. We do not claim a universally
superior model; we characterise \emph{whether} and \emph{where} cross-sectional
graph structure improves tail-risk forecasting under a disciplined protocol.

Three further points follow directly from the evidence in
Section~\ref{sec:results}. First, \emph{how} a graph is integrated matters as
much as whether one is used at all, and the effect size is large: switching
cross-sectional batch normalisation for per-node layer normalisation, and
uniform mean aggregation for learned attention, together cut the GNN-vs-MLP
gap roughly in half without changing the graph, the features, or the training
objective at all. Both of these are defaults inherited from node- and
graph-classification settings, where the prediction target is comparative
across a batch rather than a single node's idiosyncratic deviation from its
peers; our results suggest financial tail-risk GNNs should treat both choices
as task-specific design decisions rather than adopt them unmodified. Second,
attention appears to substitute for, rather than complement, some of the fixes
we tried to address graph staleness and oversmoothing directly: an
input skip-connection and a volatility-triggered graph refresh both left
GAT's score unchanged, and a GAT variant additionally conditioning attention
on $|$correlation$|$ edge weight (a separate diagnostic, not
Table~\ref{tab:ablation}) underperformed plain (unweighted) GAT. This suggests
attention is already
learning to down-weight the specific edges that a hand-designed fix would
target, making bespoke staleness/oversmoothing corrections partially redundant
once the aggregation rule itself can adapt. Third, the walk-forward
decomposition shows the graph's remaining cost is not uniform but concentrated
in regime-shift periods, and attention dropout, tuned only on the primary
evaluation's validation split, did not transfer cleanly to this very
different, expanding-window training regime: the 2022 fold's gap, if
anything, widened, while the other three folds moved in mixed directions.
This is consistent with a genuinely distinct mechanism (graph staleness
relative to the forecasting horizon) that attention alone cannot fix, since
attention can only reweight the edges it is given, not correct a graph that
is stale by construction, and that a regulariser tuned for one operating
point and split is not guaranteed to help a materially different one. A graph
refreshed at
the same or higher frequency than the forecast horizon, or one built from a
shorter, more reactive correlation window, remains a natural next step; our
own hyperparameter sensitivity study (Table~\ref{tab:sens}) did not find a
configuration that fully closes the gap, though its best cell comes close to
matching the no-graph MLP, and a wider search or an adaptive-window graph
construction remains open.

A fourth point, reported honestly in both directions because it complicates
rather than simply supports our headline result: every learned model
over-breaches VaR by roughly $2\times$ the 5\% target (Table~\ref{tab:main}),
which is consistent with a single, simple cause: the raw (VaR, ES)
\emph{magnitudes} are systematically too small, not merely mis-ranked. We test
this with the cheapest possible fix: a single scalar $k$, applied identically
to VaR and ES so $\text{ES}\ge\text{VaR}$ is preserved exactly, selected by
grid search minimising validation FZ0 and confirmed exactly once on test, the
same discipline as the ensemble weight above. The effect is large: rescaled
by $k=\recalGnnK$, the GNN's test FZ0 improves from $\gnnFZ$ to $\recalGnnFZ$,
now numerically \emph{better} than GARCH(1,1)'s $\garchFZ$ though not
significantly so ($p$\recalGnnVsGarchP, versus $p$\dmGarchP unscaled), and its
Kupiec pass rate rises from $23\%$ to $\recalGnnKupiec\%$. Rescaled by
$k=\recalMlpK$, the MLP reaches $\recalMlpFZ$ and \emph{significantly beats}
GARCH ($p$\recalMlpVsGarchP), as does the rescaled GNN+MLP ensemble
($\recalBlendFZ$, $p$\recalBlendVsGarchP), the strongest result against
GARCH anywhere in the paper. But this scalar correction does not extend the
graph's case against the no-graph model: at raw/unscaled magnitude the two are
statistically tied ($p$\dmMlpP), yet rescaling \emph{introduces} a significant
GNN-vs-MLP gap ($p$\recalGnnVsMlpP) that was not there before correcting for
scale, because the same correction that fixes each model's shared
under-confidence leaves the GNN's larger residual dispersion around the
correct scale exactly where it was, and that dispersion is what the DM test
now picks up once the shared scale bias is no longer masking it. We take this
as evidence that the (now closed) GNN-vs-MLP gap and the GNN-vs-GARCH gap have
different sources: the latter is substantially a magnitude/scale bias, fixable
with a single validation-selected number and shared across every learned model
regardless of architecture; the residual GNN-vs-MLP \emph{variance} is not,
and is better addressed by the architecture and ensemble results above, which
reduce or extract it directly rather than by rescaling.

A further limitation beyond survivorship bias: our universe and evaluation
window are US-only (2015--2024), so we cannot speak to whether cross-sectional
graph structure fares differently in markets with different microstructure,
liquidity, or crisis dynamics (e.g.\ emerging markets, where correlation
regimes are known to shift even faster). We also train and evaluate a single
graph-construction rule (sector plus thresholded correlation, monthly
refresh) as our primary specification; the construction and sensitivity
ablations bound how much this specific choice matters, but do not rule out
graph designs outside the grid we tested.

\section{Conclusion}
We use a disciplined, forward-looking evaluation protocol throughout:
strictly consistent joint (VaR, ES) scoring, correctly specified backtests,
Diebold--Mariano significance over multiple seeds, an architecture ablation
that selects the aggregation rule and its regularisation on evidence rather
than convention, and a no-graph ablation that isolates the graph's marginal
contribution. Under this protocol, a correctly regularised graph-attention
(GAT) model closes the gap to an
otherwise-identical no-graph MLP entirely at the primary 5-day 95\% tail level
for this universe of US equities ($p$\dmMlpP) and opens a significant lead
over plain GCN ($p$\dmGcnP), though it remains significantly behind
univariate GARCH, GJR-GARCH and historical simulation. This is a materially
different result than a naive GraphSAGE implementation would suggest:
switching from batch- to layer-normalisation, from uniform mean aggregation to
learned attention, and finally regularising that attention with dropout on its
coefficients, each measurably closed the GNN-vs-MLP gap in turn without
touching the graph, features or objective. The graph-construction ablation
shows edge redundancy, not graph information content, drove much of the
earlier gap: sector-only and correlation-only edges score similarly to the
combined graph once attention can down-weight uninformative edges, where
uniform aggregation could not. The rolling-origin walk-forward evaluation
shows the GNN's remaining weakness is concentrated in regime-shift periods
rather than spread uniformly across years, consistent with a graph that is,
by construction, always at least as stale as its refresh interval, a failure
mode attention can only reweight around, not correct, since GARCH's
per-series conditional-volatility recursion adapts day-by-day without
reference to any graph at all. Together this points to three distinct,
complementary mechanisms rather than a blanket verdict on graph learning for
tail risk: \emph{how} neighbour information is aggregated and how strongly its
attention is regularised (both fixed by architecture, and jointly the biggest
lever we found), and \emph{when} the graph itself is refreshed relative to the
forecasting horizon (not fixed by architecture, and the most likely source of
what gap remains). Reported honestly, that architecture and its
attention-dropout setting were chosen on primary-level validation data alone
and do not transfer uniformly: at the stricter 97.5\% (Basel/FRTB) ES level
the standalone GNN regresses significantly behind the no-graph MLP, even as it
matches or beats it at the primary level. Beyond the architecture itself, a
validation-selected ensemble of the GNN with the no-graph MLP (no additional
training, a single grid-searched weight confirmed exactly once on held-out
data at each tail level) significantly outperforms both of its components at
both tail levels, and beats GARCH(1,1) outright at 97.5\%, evidence that
ensembling extracts the graph's value more reliably than the standalone
architecture does across operating points. A separately validated scalar
recalibration shows the GNN's remaining shortfall against GARCH is
substantially a shared, fixable magnitude bias rather than a graph-specific
one: rescaled, the GNN alone statistically ties GARCH and the rescaled
ensemble significantly beats it outright, the strongest evidence in the paper
that the graph contributes real, extractable information rather than merely
underperforming less. We report this as a genuine contribution either way it
is read: a rigorous, iterated protocol for testing this question, concrete
evidence for which architectural choices matter and by how much, and a
precise characterisation of where a graph-learning approach still falls short
of simpler baselines on a decision-relevant, regulator-grade risk
measure, and where, once correctly regularised, ensembled, or recalibrated,
it no longer does.

\bibliographystyle{splncs04}
\bibliography{refs}
\end{document}
